\documentclass[11pt]{article}
\usepackage[margin=1in]{geometry}
\usepackage{amsmath,amssymb,amsthm}
\newtheorem{theorem}{Theorem}
\title{A disproof of Conjecture 00000000753}
\author{earthking11}
\date{October 6, 2026}
\begin{document}
\maketitle
\section*{Counterexample at index zero}
The conjecture proposes the exact finite count
\[
 p^{i(i+1)/2}
\]
for polynomials whose coefficients satisfy $|a_i|_p\le p^{-i}$ termwise.

\begin{theorem}
The formula fails already at $i=0$.
\end{theorem}
\begin{proof}
At $i=0$ the condition is $|a_0|_p\le1$.  Both constant polynomials
$P_0(X)=0$ and $P_1(X)=1$ satisfy it, since $|0|_p=0$ and $|1|_p=1$.
They are distinct, so there are at least two admissible polynomials.  The
displayed formula gives
\[
 p^{0(0+1)/2}=p^0=1,
\]
a contradiction.
\end{proof}

Moreover, if coefficients range in $\mathbb Q_p$ as the statement suggests,
the unit ball supplies infinitely many constant polynomials.  A finite count
would require an unstated quotient or finite coefficient alphabet.

\section*{Lean 4 certificate}
The project in \texttt{lean4/} represents constant polynomials by their
integer coefficient, records the exact index-zero integrality condition,
constructs $0$ and $1$, proves they are distinct, and evaluates the asserted
formula to $1$.  No \texttt{sorry}, \texttt{native\_decide}, or added axiom is
present.
\end{document}
